% Note for ChiPT-3 Appendix C, to follow Eq. (C25) and the paragraph after it.
% Bib entries for the new keys are in note_piNN_ohta.bib.
% Numbers: scratchpad c24_cvc.py on ChiPT_UFO_V14, vector current only,
% physical masses and form factors, W = 1.23-1.35 GeV, Q^2 = 0.1-1 GeV^2.

\paragraph*{The $\pi NN$ form factor off shell.}
In the $s$- and $u$-channel nucleon poles and in the contact term the pion at the
$\pi NN$ vertex is the external one, so $\widetilde F$ is evaluated at
$k_\pi^2 = m_\pi^2$ and reduces to its normalisation. Only the pion-in-flight
diagram has an off-shell pion at that vertex. Its momentum is $k_\pi - q$, and
its virtuality is $t = (k_\pi - q)^2$. Writing the strong vertex as
$\Gamma(k) = \widetilde F(k^2)\,\slashed{k}\gamma_5$, with
$\widetilde F(m_\pi^2) = 1$ and the coupling $g_A/2f_\pi$ and the isospin factor
stripped off, the pion-in-flight current is
\begin{equation}
  M^\mu_{\rm PF} = C\,\widetilde F(t)\,
  \frac{(2k_\pi - q)^\mu}{t - m_\pi^2}\,
  \bar u(k_2)\,(\slashed{k}_\pi - \slashed{q})\gamma_5\, u(k_1) ,
  \label{eq:MPF}
\end{equation}
where $C$ collects the couplings, the pion charge and the common $F_1^V(Q^2)$
of the $V\pi\pi$ vertex. Because $q\cdot(2k_\pi - q) = m_\pi^2 - t$, its
divergence is
\begin{equation}
  q_\mu M^\mu_{\rm PF} = -C\,\widetilde F(t)\,
  \bar u(k_2)\,(\slashed{k}_\pi - \slashed{q})\gamma_5\, u(k_1) .
\end{equation}
The nucleon poles and the Kroll--Ruderman term cancel this divergence only when
$\widetilde F(t) = 1$. Leaving $\widetilde F(t)$ on this vertex alone therefore
violates CVC by a term proportional to $1 - \widetilde F(t)$. In the model this
violation is between 3\% and 84\% of the vector current over
$W = 1.23$--$1.35$~GeV and $Q^2 = 0.1$--$1~\text{GeV}^2$, where $\widetilde F(t)$
falls from $0.79$ to $0.14$.

A momentum-dependent vertex is non-local, and gauging it consistently
introduces a $V\pi NN$ contact current. Its longitudinal part is fixed by the
Ward--Takahashi identity
$q_\mu M^\mu_{\rm c} = C\,[\Gamma(k_\pi) - \Gamma(k_\pi - q)]$. On shell,
$\Gamma(k_\pi) = \slashed{k}_\pi\gamma_5$, so the right-hand side splits into
\begin{equation}
  \Gamma(k_\pi) - \Gamma(k_\pi - q)
  = \slashed{q}\gamma_5
  + \bigl[1 - \widetilde F(t)\bigr](\slashed{k}_\pi - \slashed{q})\gamma_5 .
\end{equation}
The first term is the Kroll--Ruderman current $C\gamma^\mu\gamma_5$, which keeps
its point-like strength. The second is new. Minimal substitution in
$\widetilde F(-\partial^2)$, the prescription of Ohta~\cite{Ohta:1989}, gives
it the form
\begin{equation}
  M^\mu_{\rm O} = C\,\bigl[1 - \widetilde F(t)\bigr]\,
  \frac{(2k_\pi - q)^\mu}{t - m_\pi^2}\,
  \bar u(k_2)\,(\slashed{k}_\pi - \slashed{q})\gamma_5\, u(k_1) .
  \label{eq:MOhta}
\end{equation}
Since $\widetilde F(m_\pi^2) = 1$, the factor $1 - \widetilde F(t)$ vanishes at
the pion pole, so Eq.~\eqref{eq:MOhta} is a contact term with no pole. Added to
Eq.~\eqref{eq:MPF}, it cancels the form factor identically:
\begin{equation}
  M^\mu_{\rm PF} + M^\mu_{\rm O}
  = C\,\frac{(2k_\pi - q)^\mu}{t - m_\pi^2}\,
  \bar u(k_2)\,(\slashed{k}_\pi - \slashed{q})\gamma_5\, u(k_1) .
\end{equation}
The pion-in-flight diagram together with its gauge partner is therefore
independent of $\widetilde F$~\cite{Workman:1992}. This is why our vertices
carry no strong form factor on that diagram. Including $\widetilde F$ together
with Eq.~\eqref{eq:MOhta} reproduces our amplitude to machine precision, and
restores the residual of the full vector current to the value without
$\widetilde F$.

Three consequences follow.
\begin{itemize}
\item In single-pion production $\widetilde F$ is not observable in the vector
  current. The Ohta construction fixes only the longitudinal part of the
  contact, so any further transverse piece would be new physics with its own
  parameter, which we do not introduce.
\item The axial current has no pion-in-flight diagram, since the $A\pi\pi$
  matrix element vanishes by parity. There, $\widetilde F$ enters through the
  pion-pole terms instead: through $F_P$ in Eq.~(C15) and through $C_6^A$ in
  Eq.~(C44), where PCAC and pion-pole dominance fix it.
\item The cancellation holds diagram by diagram only because one hadron leg is
  off shell. With two or more pions in the final state, internal pions and
  nucleons are off shell in several places at once. The contact of
  Eq.~\eqref{eq:MOhta} must then be kept as an operator, and the more general
  prescription of Haberzettl~\cite{Haberzettl:1997,Haberzettl:1998} becomes
  relevant. With the crossing-symmetric choice of Davidson and
  Workman~\cite{Davidson:2001},
  $\hat F = 1 - (1-F_s)(1-F_u)(1-F_t)$, it is the $\hat F_B$ of
  Table~I, applied there at the amplitude level following Ref.~[34]. We leave
  its vertex-level implementation, which requires form factors with more than
  one argument in the UFO format, to future work.
\end{itemize}
